\section{QCD 真空中关联函数的线性发散}\label{sec:otherresult}

本节分析在 QCD 真空中测得的矩阵元里 Wilson 线的线性发散。我们试图检验这一发散是否普适，以及能否从这些矩阵元中成功地提取出发散质量。我们考虑三类新的矩阵元：1）曾用于提取重夸克势的大尺寸 Wilson 圈；2）准 PDF 算符的真空矩阵元；3）固定规范下 Wilson 链接待算符的真空矩阵元。
数据取自 MILC 系综，在适用的场合配以 clover 作用量。
按照上一节的结论，涂抹不改变重正化的形式，因此这里采用一步涂抹的高精度数据。既然只想检验线性发散，拟合中使用简单的函数式~(\ref{eq:logM_ss}) 即可。

\subsection{Wilson 圈}

要从格点数据提取线性发散，Wilson 圈是一个很好的选择，因为它计算容易、规范不变且精度高。文献中对它已有广泛的研究~\cite{Chen:2016fxx,Zhang:2017bzy,Musch:2010ka,Green:2017xeu,Chen:2017gck}，特别是与重夸克势的联系。

\begin{figure}[tbp]
\centering
\includegraphics[width=8cm]{images/fitting_Wilson_Loop_hyp1.pdf}
\includegraphics[width=8cm]{images/ert_vs_r+t_Wilson_Loop_hyp1.pdf}
\caption{HYP 涂抹情形下 Wilson 圈线性发散的检验。
其斜率与上一节强子矩阵元的值大体一致。}
\label{fig:LP_test}
\end{figure}

在我们的分析中，我们不考虑重夸克势的解释，而是简单地把 Wilson 圈当作一个矩阵元，通过查看矩阵元的对数来拟合其 $1/a$ 发散，如图~\ref{fig:LP_test} 所示。取 $\Lambda_{\rm QCD}=0.39$ GeV（上一节拟合所偏好的数值）时，格点间距依赖已被我们的公式很好地拟合出来。在分离出线性发散系数之后，我们把它画作 $r+t$ 的函数，其斜率的一半即给出发散系数。我们得到的数值约为 2.5 GeV$^{-1}$fm$^{-1}$（无量纲时为 0.49）。该数值与上一节得到的（图~\ref{fig:e(z)_hyp}）非常相近，
表明线性发散可以由 Wilson 圈很好地提取出来。

\subsection{准 PDF 算符的真空矩阵元}

\begin{figure}[tbp]
\centering
\includegraphics[width=8cm]{images/fitting_VEV_hyp1_Oa.pdf}
\includegraphics[width=8cm]{images/ez_vs_z_VEV_hyp1_Oa.pdf}
\caption{准 PDF 算符真空矩阵元的线性发散检验。矩阵元被强烈压低，不存在领头扭度贡献。}
\label{fig:vev12}
\end{figure}

这里我们考虑准 PDF 算符的真空矩阵元，文献~\cite{Braun:2018brg} 曾建议将其作为重正化因子的候选。

可以考虑 $O_{\gamma_t}(t)$ 的真空期望值（VEV）。作为文献~\cite{Braun:2018brg} 提出的规范不变选择，$\langle O_{\Gamma}(t)\rangle$
可通过如下随机估计获得：
\bea\label{eq:vev1}
&&\langle O_{\Gamma}(t)\rangle =\frac{1}{L_s^3}\langle \sum_{\vec{x}}\textrm{Tr}[U(\vec{x},0;\vec{x},t)\Gamma S_w(\vec{x},t)]\rangle\nonumber\\
&=&\frac{1}{L_s^3}\langle \sum_{\vec{x}}\textrm{Tr}[U(\vec{x},0;\vec{x},t)\Gamma S(\vec{x},t;\vec{x},0)]\rangle\nonumber\\
&&+\frac{1}{L_s^3}\sum_{\vec{x},\vec{y},\vec{y}\neq\vec{x}}\langle \textrm{Tr}[U(\vec{x},0;\vec{x},t)\Gamma S(\vec{x},t;\vec{y},0)]\rangle
\eea
其中第二行右端的第二项不是规范不变的，因而为零；$S_w(\vec{x},t)=\sum_{\vec{y}}S(\vec{x},t;\vec{y},0)$ 是不经规范固定的墙源夸克传播子。这样的提案只适用于 $\Gamma=\gamma_t$ 或 ${\cal I}$ 的非零情形，我们将把它应用于 $O_{\gamma_t}$。

然而，利用算符乘积展开容易看出，真空中 $O_{\gamma_t}$ 矩阵元在小 $t$ 处为零。
因此该矩阵元容易受到大的 O($a$) 修正的影响；只有在非常大的 $t$ 处才可能看到正确的斜率。

PDF 算符真空矩阵元线性发散的检验结果示于图~\ref{fig:vev12}。
大 $t$ 处的斜率看起来与 Wilson 圈情形一致，但只是在毫无意义的大误差范围内一致。

\subsection{Landau 规范固定的 Wilson 链接}

对规范依赖的矩阵元，我们考虑最简单的选择：Landau 规范下的 Wilson 线。

\begin{figure}[tbp]
\centering
\includegraphics[width=8cm]{images/fitting_Wilson_Link_hyp1.pdf}
\includegraphics[width=8cm]{images/ez_vs_z_Wilson_Link_hyp1.pdf}
\caption{HYP 涂抹情形下 Landau 规范固定 Wilson 链接的线性发散检验。小 $z$ 处可用微扰论计算，线性发散表现良好；大 $z$ 处该矩阵元在此规范下没有转移矩阵的诠释。}
\label{fig:link_test}
\end{figure}

我们检验的结果示于图~\ref{fig:link_test}。
在微扰论适用的小 $z$ 区，斜率与微扰论的预言大体相符；大 $z$ 处该矩阵元在此规范下没有转移矩阵的诠释。

\begin{figure}[tbp]
\centering
\includegraphics[width=8cm]{images/ez_vs_z_Comp_Oa.pdf}
\caption{不同真空矩阵元中线性和发散的比较。}
\label{fig:com_test}
\end{figure}

最后，图~\ref{fig:com_test} 比较 Wilson 圈、VEV 与 Wilson 链接三者的线性发散项。小 $z$ 处 Wilson 圈与 Wilson 链接的线性发散项相互一致，而大 $z$ 处则强烈分歧。VEV 的线性发散项与其他二者不同，因为大的离散化误差使其提取非常不精确。
